\chapter{Tiny210
开发者指南}\label{tiny210-ux5f00ux53d1ux8005ux6307ux5357}

本文档集面向在 FriendlyARM Tiny210（Samsung S5PV210、ARM
Cortex-A8）上开发 xv6 的人员。默认目标是 Tiny210；同一端口也可通过
\passthrough{\lstinline!S5PV210\_BOARD=x210!} 构建九鼎 X210，但两块板的
DRAM、串口和机器号不同。

\begin{quote}
当前状态：S5PV210 实板路径与 QEMU 替身路径共用 ARM32 内核框架。QEMU 使用
\passthrough{\lstinline!cubieboard!}（Allwinner A10）验证通用
ARM32、initcall、platform bus、 进程和文件系统逻辑，并不能验证 S5PV210
的 VIC、PWM、FIMD、DDR 时序和引脚。 修改寄存器驱动后仍需 Tiny210
实板验证。
\end{quote}

\section{阅读顺序}\label{ux9605ux8bfbux987aux5e8f}

\begin{enumerate}
\def\labelenumi{\arabic{enumi}.}
\tightlist
\item
  \href{boot.md}{启动与内存初始化}：从 iROM/BL1 到第一个用户进程。
\item
  \href{interrupts.md}{异常、中断与系统时钟}：异常向量、VIC、PWM Timer 4
  和调度 tick。
\item
  \href{lcd.md}{LCD/FIMD}：面板时序、FIMD DMA scanout、framebuffer 与
  fbcon。
\item
  \href{wifi.md}{MT7601U Wi-Fi}：EHCI 枚举、固件、SoftMAC、WPA2、NAPI 和
  DHCP。
\item
  \href{filesystem.md}{文件系统}：内嵌 xv6fs、bcache、redo log 和
  \passthrough{\lstinline!/init!}。
\item
  \href{drivers.md}{驱动模型与板载按键}：platform
  device/driver、resource、initcall、 K1-\/-K8 输入子系统和
  \passthrough{\lstinline!event0!}。
\item
  \href{bringup-notes.md}{上板调试记录}：串口乱码、EHCI 欠载、iROM 读第
  0 扇区卡死、BL1 入口偏移与对齐错误等问题的现象、原因和解决办法。
\end{enumerate}

\section{代码地图}\label{ux4ee3ux7801ux5730ux56fe}

\begin{lstlisting}
ports/tiny210/
|-- Makefile             Tiny210/X210 构建入口
|-- mktiny210.py         BL1、uImage、整卡镜像生成与校验
`-- test_shell.py        QEMU shell 回归测试

kernel/arch/arm/
|-- vectors.S            ARMv7 异常向量和 trapframe 保存/恢复
|-- irq.c                request_irq() 与通用 IRQ 分派
|-- time.c               clocksource、ticks、调度 tick
|-- vm.c                 ARMv7 short-descriptor 页表和 MMU
`-- initcall.c           按 level/order 执行内置驱动

kernel/arch/arm/mach-s5pv210/
|-- bl1/                 自带第一阶段引导程序
|-- entry.S              内核物理入口、各模式栈、VBAR
|-- setup.c              x210_main() 和板级启动
|-- devices.c            S5PV210 platform device/resource 表
|-- board.h              Tiny210/X210 板级差异
|-- memlayout.h          DRAM、用户空间和 framebuffer 布局
`-- atag.c               ARM Linux ATAG 解析

kernel/drivers/
|-- clk/clk-s5pv210.c
|-- irqchip/irq-vic.c
|-- clocksource/samsung_pwm_timer.c
|-- tty/serial/samsung.c
|-- watchdog/s3c2410_wdt.c
|-- video/fbdev/s3c-fb.c
|-- video/fbdev/fbcon.c
|-- input/keyboard/s5pv210-keys.c
|-- usb/host/ehci-hcd.c
|-- usb/core/usb.c
|-- net/wireless/mt7601u.c
`-- block/ramdisk_embed.c
\end{lstlisting}

\section{快速构建}\label{ux5febux901fux6784ux5efa}

在仓库根目录执行：

\begin{lstlisting}[language=sh]
make tiny210
\end{lstlisting}

生成整卡镜像并写入已挂载的 SD 卡（BL1 到第 1 扇区、uImage 到第 49
扇区，MBR 不动，写后读回校验）：

\begin{lstlisting}[language=sh]
make tiny210-sdimg
\end{lstlisting}

等价的显式 Tiny210 构建是：

\begin{lstlisting}[language=sh]
make -f ports/tiny210/Makefile \
  S5PV210_BOARD=tiny210 \
  S5PV210_LCD=s70
\end{lstlisting}

主要产物位于 \passthrough{\lstinline!build/tiny210/!}：

{\def\LTcaptype{none} % do not increment counter
\begin{longtable}[]{@{}p{0.34\linewidth}p{0.60\linewidth}@{}}
\toprule\noalign{}
产物 & 作用 \\
\midrule\noalign{}
\endhead
\bottomrule\noalign{}
\endlastfoot
\passthrough{\lstinline!tiny210-bl1.bin!} & 16 KiB、带 iROM 头的
BL1，写入 SD 扇区 1 \\
\passthrough{\lstinline!xv6-tiny210.uImage!} & legacy uImage；写在
第 49 扇区，是 BL1 唯一的启动内核 \\
\passthrough{\lstinline!tiny210-sd.img!} & BL1@1 与 uImage@49
已放到固定扇区的整卡原始镜像 \\
\passthrough{\lstinline!kernel.elf!} /
\passthrough{\lstinline!kernel.map!} & 带符号调试文件和链接映射 \\
\passthrough{\lstinline!fs.img!} & 4 MiB xv6fs；构建时会被嵌入内核
\passthrough{\lstinline!.ramdisk!} \\
\end{longtable}
}

构建与 QEMU 回归测试：

\begin{lstlisting}[language=sh]
make tiny210-check
make tiny210-qemu
\end{lstlisting}

更换 \passthrough{\lstinline!S5PV210\_BOARD!}、LCD 型号或
\passthrough{\lstinline!BUILD!} 配置时先清理，避免复用旧目标文件：

\begin{lstlisting}[language=sh]
make tiny210-clean
\end{lstlisting}

\section{默认硬件参数}\label{ux9ed8ux8ba4ux786cux4ef6ux53c2ux6570}

{\def\LTcaptype{none} % do not increment counter
\begin{longtable}[]{@{}p{0.34\linewidth}p{0.60\linewidth}@{}}
\toprule\noalign{}
项目 & Tiny210 默认值 \\
\midrule\noalign{}
\endhead
\bottomrule\noalign{}
\endlastfoot
CPU & Samsung S5PV210 / Cortex-A8，单核 \\
DRAM & 512 MiB，物理起始 \passthrough{\lstinline!0x20000000!} \\
内核入口 & \passthrough{\lstinline!0x20008000!} \\
ATAG & \passthrough{\lstinline!0x20000100!} \\
机器号 & 3466（MINI210） \\
串口 & UART0 / COM0，115200 8N1 \\
IRQ 控制器 & 4 个 PL192 VIC，IRQ 0..127 \\
系统 tick & PWM Timer 4，IRQ 25 \\
LCD & S70 800×480，XRGB8888（可编译选择其他面板） \\
根文件系统 & 优先 FAT32 中的
\passthrough{\lstinline!FS.IMG!}，找不到时回退内嵌 xv6fs RAM 镜像 \\
\end{longtable}
}

\section{启动全景}\label{ux542fux52a8ux5168ux666f}

\begin{figure}[htbp]
\centering
\resizebox{\linewidth}{!}{\begin{tikzpicture}[node distance=5mm]
  \node[hbox] (rom) {S5PV210 iROM\\BL0};
  \node[box,right=of rom] (bl1) {BL1 / iRAM\\时钟、UART、DDR};
  \node[box,right=of bl1] (mbr) {LBA 0 / MBR\\定位 FAT32};
  \node[box,right=of mbr] (img) {根目录文件\\加载 uImage};
  \node[box,right=of img] (entry) {\texttt{\_entry}\\ARM32 环境};
  \node[box,right=of entry] (main) {\texttt{x210\_main()}\\MMU、设备、进程};
  \node[hbox,right=of main] (user) {\texttt{/init}\\\texttt{/bin/sh}};
  \draw[arr] (rom)--(bl1); \draw[arr] (bl1)--(mbr); \draw[arr] (mbr)--(img); \draw[arr] (img)--(entry);
  \draw[arr] (entry)--(main); \draw[arr] (main)--(user);
\end{tikzpicture}}
\caption{Tiny210 从上电到第一个用户 shell 的主路径}
\label{fig:tiny-overview}
\end{figure}


\section{当前实现边界}\label{ux5f53ux524dux5b9eux73b0ux8fb9ux754c}

\begin{itemize}
\tightlist
\item
  内核使用 ARMv7 short-descriptor 页表和恒等映射；cache 在当前 bring-up
  中关闭。
\item
  只启用单 CPU；没有 SMP/GIC 路径。
\item
  UART 输入不是 UART IRQ 驱动，而是每次 PWM tick 调用
  \passthrough{\lstinline!uartintr()!} 轮询。
\item
  K1-\/-K8 使用 GPH2/GPH3 外部中断和 20 ms 延迟消抖，通过
  \passthrough{\lstinline!/dev/input/event0!} 输出 Linux 风格
  \passthrough{\lstinline!EV\_KEY!}/\passthrough{\lstinline!EV\_SYN!}
  事件。
\item
  LCD 只有 FIMD Window 0、单 framebuffer、32 bpp；没有 VSync
  IRQ、双缓冲或用户态 mmap。
\item
  实板构建已包含 SDHCI、FAT32、VFS 和 ext2；QEMU 构建仍使用内嵌 RAM
  镜像。
\item
  USB-A/EHCI 和 MT7601U 已接入，但 EHCI 尚无运行时热插拔，QEMU
  也无法验证。
\item
  \passthrough{\lstinline!DEVICE\_MAX!} 和字符设备表等仍是静态上限，不是
  Linux 的动态对象模型。
\end{itemize}

\section{实板调试入口}\label{ux5b9eux677fux8c03ux8bd5ux5165ux53e3}

\begin{itemize}
\tightlist
\item
  串口先看 BL1 是否打印；完全无输出时优先检查启动拨码、COM0、115200
  8N1。
\item
  BL1
  可单独构建串口测试：\passthrough{\lstinline!make -f ports/tiny210/Makefile uarttest!}。
\item
  内核加载前失败依次看
  \passthrough{\lstinline!tiny210-bl1.bin!}、第 49 扇区的
  \passthrough{\lstinline!xv6-tiny210.uImage!} 和 uImage CRC；
  根文件系统的 FAT32/raw 探测错误属于内核侧；内核入口失败看
  \passthrough{\lstinline!\_entry!}。
\item
  进入内核后按 \passthrough{\lstinline!initcall:!} 日志判断失败属于
  clock、VIC、timer、serial 还是 LCD probe。
\item
  LCD 无显示但串口正常时，先选择正确
  \passthrough{\lstinline!S5PV210\_LCD!}，再核对 RGB 排线、供电和背光。
\item
  完整烧录步骤与风险提示见
  \href{../../../ports/tiny210/README.md}{ports/tiny210/README.md}。
\end{itemize}
